\chapter*{Заключение}
\addcontentsline{toc}{chapter}{Заключение}

В ходе выполнения производственной практики была разработана платформа выявления аномалий сетевой активности на основе потокового анализа сетевых данных.

\textbf{Основные результаты работы:}

\begin{enumerate}
    \item Проведён анализ существующих подходов к мониторингу и обнаружению аномалий сетевой активности. Рассмотрены сигнатурные методы (Snort, Suricata, Zeek), статистические методы и методы машинного обучения (Isolation Forest, LOF, One-Class SVM). На основе анализа сформулированы 8 функциональных и 5 нефункциональных требований к разрабатываемой платформе.

    \item Спроектирована модульная архитектура платформы, включающая компоненты сбора данных, потоковой обработки, обнаружения аномалий и визуализации. Архитектура обеспечивает слабое связывание модулей, возможность их независимого тестирования и расширения.

    \item Реализован модуль сбора сетевого трафика, поддерживающий live-захват через Scapy, загрузку PCAP-файлов и генерацию синтетических данных с четырьмя типами атак: DDoS, port scan, bruteforce, data exfiltration.

    \item Разработан модуль потоковой обработки, выполняющий агрегацию пакетов во временные окна с вычислением 23 статистических признаков, включая энтропийные метрики для оценки разнообразия IP-адресов и портов.

    \item Реализованы и протестированы три алгоритма машинного обучения (Isolation Forest, LOF, One-Class SVM) и детектор на эвристических правилах (5 правил). Проведено сравнение методов: лучший результат по F1-мере показал LOF (0.333), все методы обеспечили полноту обнаружения Recall = 1.0.

    \item Разработан модуль визуализации, автоматически генерирующий 8 типов графиков для анализа результатов: временные ряды с аномалиями, confusion matrix, ROC-кривая, PR-кривая, распределение оценок, корреляция признаков, сравнение методов.

    \item Написаны 19 автоматизированных тестов (8 для модуля сбора, 11 для детекторов), все успешно проходят. Проведено интеграционное тестирование полного цикла обработки данных.
\end{enumerate}

\textbf{Направления дальнейшего развития:}

\begin{enumerate}
    \item Интеграция с системами потоковой передачи данных (Apache Kafka, RabbitMQ) для обработки больших объёмов трафика в реальном времени.
    \item Добавление нейросетевых методов обнаружения (Autoencoder, LSTM) для повышения качества детектирования сложных аномалий.
    \item Поддержка дополнительных сетевых протоколов (HTTP, DNS, DHCP) для более глубокого анализа.
    \item Разработка веб-интерфейса для мониторинга в реальном времени.
    \item Внедрение механизма активного обучения (active learning) для снижения доли ложных срабатываний.
\end{enumerate}

Разработанная платформа может быть использована как в учебных целях для изучения методов обнаружения сетевых аномалий, так и в качестве основы для создания корпоративных систем мониторинга информационной безопасности.
